\subsection{Classi coperte e limiti della strategia}
Riferimenti: \path{backend/BugBoard26.Api/Services/IssueDomainService.cs} e
\path{backend/BugBoard26.Tests/IssueDomainServiceTests.cs}. La strategia usa
classi di equivalenza e regole decisionali; non costituisce una suite completa
SECT/R-WECT e non dimostra copertura strutturale completa.

\textbf{SearchIssues.} Le keyword presenti sono \texttt{login}, \texttt{API}
e \texttt{grafica}: coprono titolo, descrizione e assenza di risultati.
Il test asserisce il conteggio, non gli ID restituiti. Keyword nulla, vuota,
di soli spazi, trimming, differenze di maiuscole, elenco vuoto e molteplici
corrispondenze non hanno casi dedicati in questo file. In particolare, la
keyword API e' identica al testo della fixture: non prova da sola il confronto
case-insensitive del metodo. Le prove PostgreSQL della ricerca API restano
distinte da quelle dell'helper in memoria.

\textbf{CanChangeIssueStatus.} Le classi coperte separano ADMIN attivo,
USER assegnatario/non assegnatario, READONLY assegnatario e USER inattivo.
L'oracolo e' il booleano atteso. Non sono coperte tutte le combinazioni:
per esempio ADMIN inattivo e USER attivo con assegnatario nullo non hanno
casi unitari dedicati. Date, transizioni effettive e salvataggi appartengono
al controller e richiedono test HTTP/persistenza, non questo metodo puro.

\textbf{SuggestAssignee.} Le Theory usano carichi (0,2) e (3,1) e verificano
l'ID del vincitore; le fixture escludono un READONLY e issue Resolved/Closed.
Il Fact usa soltanto READONLY attivo e USER inattivo e attende \texttt{null}.
Parita' per ID, ADMIN candidato, InProgress, archiviati ed elenco vuoto non
hanno casi unitari dedicati nel file. Il codice implementa queste regole e
sono presenti scenari HTTP aggiuntivi; non si trasforma la loro presenza
in copertura unitaria misurata. La suite del 19 ottobre esegue anche questi
test HTTP: 166 casi complessivi superati, 11 di dominio e 155 di integrazione.
Il report e il TRX sono in \path{verifica-backend/}; non si attribuiscono
alle sole prove unitarie i risultati ottenuti attraverso l'API.
